\section{Exact HLP levels and the smooth packet local replacement}
\label{sec:hlp-local}

Fix once and for all a spectral right edge
\[
\Re s=1+\delta,\qquad \delta>0.
\]
All Dirichlet expansions in this section are first made on this line. Absolute
convergence gives $P(s)=O_\delta(1)$ and $Q(s)=O_\delta(1)$, while
$f(s)=\frac12L+O_\delta(1)$ on $t\asymp T$. Hence
\[
\left|\frac{P(s)}{f(s)}\right|\ll_\delta L^{-1},
\]
The coefficient organization is the one appearing in the local-extrema analysis
of Hughes--Lugmayer--Pearce-Crump~\cite{HLP}, but no RH-dependent estimate from
that paper is used here.  Indeed, on the present fixed safe line the exact
identity
\begin{equation}
\boxed{
\frac{Q(s)}{f(s)-P(s)}
=\sum_{k\ge1}f(s)^{-k}P(s)^{k-1}Q(s)
}
\label{eq:safe-line-HLP-expansion}
\end{equation}
follows directly from the geometric series, and converges absolutely and
uniformly because $|P/f|\ll_\delta L^{-1}$.  Thus the levelwise expansion used
below is unconditional before any contour movement.  Since $f'(s)=O(T^{-1})$,
its separated $f'$ contribution is also $O_\delta((T L)^{-1})$ on this line.
The exact multiplier at level $k$ changes only the Dirichlet coefficient
sequence. If the total convolution index is $M$, the stationary phase remains
\begin{equation}
t_*=2\pi My,
\qquad
\text{carrier}=e(-My).
\label{eq:total-carrier}
\end{equation}
Thus exact levels do not introduce a new moving support variable. Their packet
weights remain in the finite-interval Laplace/Mellin class
\eqref{eq:joint-weight}.

\subsection{Uniform factorial control of the HLP hierarchy}
For $k\ge0$ define the nonnegative Dirichlet-convolution coefficients
$\Lambda_k$ by
\[
P(s)^k=\sum_{n\ge1}\frac{\Lambda_k(n)}{n^s},
\qquad \Lambda_0(1)=1,
\]
and for $k\ge1$ let $\alpha_k$ be the coefficient sequence of
$P^{k-1}Q$.

\begin{lemma}[Derivative identity and support]
\label{lem:alpha-identity}
For every $k\ge1$,
\begin{equation}
\boxed{
\alpha_k(n)=\frac{\log n}{k}\Lambda_k(n),
\qquad
\Lambda_k(n)\ne0\Longrightarrow n\ge2^k.
}
\label{eq:alpha-identity}
\end{equation}
\end{lemma}

\begin{proof}
Differentiate the identity for $P^k$. Since $P'=-Q$,
\[
-\sum_n(\log n)\Lambda_k(n)n^{-s}
=kP^{k-1}P'
=-kP^{k-1}Q
=-k\sum_n\alpha_k(n)n^{-s}.
\]
Coefficient comparison gives the first identity. Every nonzero term of the
$k$-fold convolution defining $\Lambda_k$ is a product of integers at least
$2$, giving the support statement.
\end{proof}

\begin{proposition}[Uniform factorial square-sum majorant]
\label{prop:factorial-majorant}
There are absolute constants $A,C>0$ such that, for every $x\ge2$ and
$k\ge1$,
\begin{equation}
\boxed{
\sum_{n\le x}\alpha_k(n)^2
\le C\frac{A^k}{k!}\,
 x\bigl(\log(ex)\bigr)^{2k+1}.
}
\label{eq:factorial-majorant}
\end{equation}
\end{proposition}

\begin{proof}
Put
\[
S_k(x):=\sum_{n\le x}\Lambda_k(n)^2.
\]
The convolution identity
\[
\Lambda_k(n)=\sum_{d\mid n}\Lambda(d)\Lambda_{k-1}(n/d)
\]
and $\sum_{d\mid n}\Lambda(d)=\log n$ give by weighted Cauchy--Schwarz
\[
\Lambda_k(n)^2
\le(\log n)\sum_{d\mid n}\Lambda(d)\Lambda_{k-1}(n/d)^2.
\]
Hence
\begin{equation}
S_k(x)\le(\log x)\sum_{d\le x}\Lambda(d)S_{k-1}(x/d).
\label{eq:Sk-recurrence}
\end{equation}
Chebyshev's estimate $\psi(y)=\sum_{n\le y}\Lambda(n)\ll y$ and Stieltjes
integration by parts imply, uniformly for $0\le v\ll\log(ex)$,
\begin{equation}
\sum_{d\le x}\frac{\Lambda(d)}d
\left(\log\frac{ex}{d}\right)^v
\ll\frac{(\log(ex))^{v+1}}{v+1}.
\label{eq:weighted-Chebyshev}
\end{equation}
Indeed, apply partial summation to
$d^{-1}(\log(ex/d))^v$ and use $\psi(t)\ll t$.

Inducting in \eqref{eq:Sk-recurrence}, starting from
$S_1(x)\ll x\log(ex)$, and applying
\eqref{eq:weighted-Chebyshev} with $v=2k-3$ gives
\[
S_k(x)
\le C\frac{A^{k-1}}{(k-1)!}
 x(\log(ex))^{2k-1}.
\]
If $S_k(x)\ne0$, Lemma~\ref{lem:alpha-identity} gives $2^k\le x$, so the
required condition $v\ll\log(ex)$ is automatic. Finally
\[
\sum_{n\le x}\alpha_k(n)^2
\le\frac{(\log x)^2}{k^2}S_k(x),
\]
and enlarging $A,C$ proves \eqref{eq:factorial-majorant}.
\end{proof}

\begin{lemma}[Quadratic bound for completed direct carriers]
\label{lem:direct-carrier-square}
After the common shifted-zeta factor has been placed in the completed analysis
map and the deterministic stationary/shell normalization defining
$\mathcal K_{j,T}$ has been applied, let $c_m^{\rm dir}(p)$ denote the remaining
normalized coefficient of a direct one-$\chi$ carrier. For the ordinary
logarithmic-derivative term, the exact HLP hierarchy, and their translated
variants, uniformly for $x\le T$ and $0\le q\le2$,
\begin{equation}
\boxed{
\sum_{m\le x}
\left|\partial_p^q c_m^{\rm dir}(p)\right|^2
\ll xL^{C_{\rm dir}}.
}
\label{eq:direct-carrier-square}
\end{equation}
The same bound holds for any subset of the HLP levels, in particular
$k>K_0$, for any support restriction of these direct carriers, and uniformly
for the translated principal-part subfamilies arising from the locally
projected levels $k\le K_0$ in \eqref{eq:canonical-local-decomp}.
\end{lemma}

\begin{proof}
By Proposition~\ref{prop:trivial-outer-twist}, no nontrivial outer arithmetic
map remains. The direct terms considered here are precisely the terms for which
Lemma~\ref{lem:internal-principal-contraction} leaves the radial carrier $M$
uncontracted: the additional translated residues and the exact levels above
$K_0$ retain their direct carrier. Thus for $k\ge1$ the changing internal
multiplier is exactly $f^{-k}P^{k-1}Q$, with coefficient
$f^{-k}\alpha_k(m)$. Since $|f|\asymp L$ and $\log(ex)\ll L$ for $x\le T$,
\eqref{eq:factorial-majorant} gives
\[
\left(\sum_{m\le x}|f^{-k}\alpha_k(m)|^2\right)^{1/2}
\ll (xL)^{1/2}\frac{C^k}{\sqrt{k!}}.
\]
A fixed $p$-derivative introduces only a polynomial in $k$ and factors
$(\log m)/L=O(1)$. Minkowski's inequality and
$\sum_k C^k(1+k)^C/\sqrt{k!}<\infty$ therefore give the same bound after
summing any subset of levels. The ordinary $k=0$ coefficient is controlled by
$\sum_{m\le x}\Lambda(m)^2\le (\log x)\psi(x)\ll xL$; translated factors
$m^{O(1/L)}$ are uniformly bounded. A support restriction only sets some of the
Hilbert-valued coefficient coordinates to zero and therefore cannot increase
the square sum. For a fixed level $k$, each additional translated principal part
is obtained by a finite residue functional on a normalized local contour and
does not introduce a new carrier index. Cauchy's estimate and the fixed-factor
bounds used in \eqref{eq:uniform-Hk} give at most
$L^C C_1^k(1+k)^C$ beyond the level coefficient bound. Hence Minkowski's
inequality reduces the complete translated family to the convergent series
\[
\sum_{k\ge1}\frac{C_2^k(1+k)^C}{\sqrt{k!}}<\infty.
\]
Thus the translated principal-part contribution is uniform in $K_0$ and obeys
\eqref{eq:direct-carrier-square} after enlarging $C_{\rm dir}$.
\end{proof}

\begin{lemma}[Uniform freezing of the slow half-self factor]
\label{lem:block-freezing}
Let a smooth height block be supported on
$|t-\tau_b|\ll H=T/L^D$ on the fixed right edge $\Re s=c>1$, and let
$F_b=f(1/2+i\tau_b)$ be as in \eqref{eq:Fb-real-freeze}. Then
\begin{equation}
\boxed{
|f(c+it)-F_b|\ll L^{-D},
\qquad
|f(c+it)^{-k}-F_b^{-k}|
\ll kL^{-k-D-1}\quad(k\ge1).
}
\label{eq:block-freezing-coeff}
\end{equation}
After restricting the HLP hierarchy to $k\le K_0$ and applying the local
principal projector, there is a fixed $C_{\rm fr}$ such that the normalized
local kernels obtained from \eqref{eq:Hex-variable} and \eqref{eq:Hex} satisfy,
for the derivative orders used below,
\begin{equation}
\boxed{
\partial_\rho^r\partial_p^q
\bigl(K_{\rm var,b}-K_{\rm fr,b}\bigr)
=O\!\left(L^{-D-1+C_{\rm fr}}\right),
\qquad 0\le r,q\le2.
}
\label{eq:block-freezing-kernel}
\end{equation}
The same estimate holds for the corresponding diagonal derivative jump.
\end{lemma}

\begin{proof}
Stirling's formula on a fixed vertical line gives
$f(c+it)=\frac12\log(t/2\pi)+O(T^{-1})$ and
$\partial_t f(c+it)=O(T^{-1})$. Moreover $f'(s)=O(T^{-1})$ uniformly on the
bounded horizontal segment joining $1/2+i\tau_b$ to $c+i\tau_b$. Therefore
\[
|f(c+it)-F_b|
\ll \frac{|t-\tau_b|}{T}+T^{-1}
\ll L^{-D}.
\]
Since both $f(c+it)$ and $F_b$ have modulus
$\asymp L$, the mean-value theorem applied to $z\mapsto z^{-k}$ gives the
second estimate in \eqref{eq:block-freezing-coeff}.

The normalized self term contains the factor $2f/L$, so its freezing error is
$O(L^{-D-1})$. At HLP level $k$, every occurrence of the slow factor is
multiplied by $f^{-k}$; after a fixed number of stationary and packet
derivatives the loss is bounded by
$L^{C_{\rm fr}}(1+k)^{C_{\rm fr}}$. Proposition~\ref{prop:factorial-majorant}
therefore bounds the required sum uniformly in $K_0$, since
\[
\sum_{1\le k\le K_0}
\frac{R^k(1+k)^{C_{\rm fr}+1}}{\sqrt{k!}}
\le
\sum_{k\ge1}
\frac{R^k(1+k)^{C_{\rm fr}+1}}{\sqrt{k!}}<\infty.
\]
The finite Cauchy and curvature factors have uniform derivatives on the fixed
normalized shift polydisc. This proves
\eqref{eq:block-freezing-kernel}, including the one-sided diagonal derivatives.
\end{proof}

\begin{lemma}[High exact HLP levels remain direct sideband terms]
\label{lem:HLP-high-edge}
Let $K_0=L^{1/8}$. Keep the stationary levels $k>K_0$ in their original
pointwise-$f(t)$, absolutely convergent safe-line Dirichlet representation; no
block freezing or local Perron projector is applied to them. On a retained unit shell put
$d_j\asymp T/J_j$ and
\[
R_j^{\rm sb}:=\epsilon_Ld_j,
\]
with $\epsilon_L=L^{-\sigma}$ fixed in Section~\ref{sec:retained}. Expand the
smooth stationary amplitude in shell Fourier modes $r$. The part
$|r|\le R_j^{\rm sb}$ and its adjoint define $E_T^{\rm HLP,hi}$ with
\begin{equation}
\boxed{
\rank E_T^{\rm HLP,hi}\ll T=o(N_T).
}
\label{eq:HLP-high-edge}
\end{equation}
Write $R_T^{\rm HLP,hi}$ for the complementary part
$|r|>R_j^{\rm sb}$; it is disjoint from $E_T^{\rm HLP,hi}$ and is estimated
later as regular class~(vii).
\end{lemma}

\begin{proof}
The safe-line geometric expansion converges absolutely. Each level changes only
the coefficient of the total index $M$, and
Lemma~\ref{lem:hlp-radial-variable} gives carrier $e(-My)$ with
$t_*=2\pi My$. Lemma~\ref{lem:stationary-symbol} is uniform over the summed
exact HLP hierarchy, so its shell amplitude has the same Fourier-sideband decay
used in class~(vii). The low sidebands replace $M$ by the effective carrier
$M-r$ and hence lie in the fixed edge spaces of
Proposition~\ref{prop:edge-rank}. By Proposition~\ref{prop:trivial-outer-twist}
there is no additional arithmetic multiplicity, which gives
\eqref{eq:HLP-high-edge}. The complementary sidebands are exactly the
high-sideband class~(vii).
\end{proof}

\subsection{Canonical local pole projection and the smooth packet Mellin class}
\label{sec:local-projector}

The principal local projection is defined levelwise. Expand the exact
resolvent geometrically on the absolute-convergence right edge and let
$\mathcal F_{k,L}(1+u;p)$ denote one completed retained level after the finite
three-$Z$/Cauchy completion has been formed. Its common principal local singularity is the complete Laurent part at
$u=0$, including the actual multiplicity created by coalescing faces. Define
\begin{equation}
\Pi_0^{(k)}\mathcal F_{k,L}
:=\operatorname{PP}_{u=0}\mathcal F_{k,L}.
\label{eq:levelwise-PP-projector}
\end{equation}
If that completed level has additional translated one-$\chi$ poles separated
from the common local cluster, let $\Sigma_{{\rm tr},k}$ be their finite set.
Its cardinality is bounded independently of $k$: these poles arise only from
the fixed three-$Z$/Cauchy completion, while the HLP level changes the pole
multiplicity at the common point $u=0$. Put
\begin{equation}
\Pi_{\Sigma_{{\rm tr},k}}F
:=\sum_{u_0\in\Sigma_{{\rm tr},k}}
\operatorname{PP}_{u=u_0}F.
\label{eq:local-PP-projector}
\end{equation}
Then the exact level has the mutually exclusive decomposition
\begin{equation}
\boxed{
\mathcal F_{k,L}
=\Pi_0^{(k)}\mathcal F_{k,L}
+\Pi_{\Sigma_{{\rm tr},k}}\mathcal F_{k,L}
+\mathcal F_{k,L}^{\rm hol},
}
\label{eq:canonical-local-decomp}
\end{equation}
where $\mathcal F_{k,L}^{\rm hol}$ is holomorphic throughout the local
neighbourhood. Identity \eqref{eq:canonical-local-decomp} is algebraic for every
$k$, but the exact source decomposition applies its local projectors only for
$k\le K_0$. Thus the exact locally projected functional is the finite sum
\begin{equation}
\mathcal L_{\rm pr,\le K_0}^{\rm ex}[W]
:=\sum_{0\le k\le K_0}\omega_{k,L}
\Res_{u=0}\!\bigl(\mathcal F_{k,L}(1+u;p)W(u)\bigr).
\label{eq:truncated-exact-principal-functional}
\end{equation}
The levels $k>K_0$ are not projected and remain in the direct sideband
decomposition of Lemma~\ref{lem:HLP-high-edge}.

For comparison with the positive model, replace the leading Laurent factor at
each level and resum the resulting \emph{model} hierarchy over all $k$. This is
not a second extraction of the exact high-level source. Lemma~\ref{lem:leading-jet-resummation}
then produces the model cluster
\begin{equation}
\boxed{\Sigma_0^{\rm mod}=\{0,2/L\},}
\label{eq:model-principal-cluster}
\end{equation}
with residues $-2,+1$. For the resummed model define
\begin{equation}
\Pi_{\Sigma_0^{\rm mod}}F
:=\sum_{u_0\in\{0,2/L\}}\operatorname{PP}_{u=u_0}F.
\label{eq:model-PP-projector}
\end{equation}
Thus the translated model pole belongs to the \emph{model} principal
projector, while the algebraic point $u_*(F_b)$ of
Lemma~\ref{lem:principal-pole-cluster} serves only to identify the same
translated scale. The levelwise decomposition therefore separates the common
principal cluster, the additional translated poles, and the holomorphic
remainder into mutually exclusive contributions.

\begin{lemma}[Smooth bounded-support Mellin weights]
\label{lem:smooth-mellin-class}
Let $0\le t\le C_*$ with $C_*$ fixed. Assume that for some fixed derivative
order $R_*$, large enough for the applications below,
\[
\|\partial_t^m\partial_\rho^r\partial_p^j a_{\rho,p,T}\|_{L^1(0,C_*)}
\le L^{A_*+o(1)}
\qquad (m+r+j\le R_*).
\]
Define
\begin{equation}
W_{L,\rho,p}(u)
:=\int_0^{C_*}a_{\rho,p,T}(t)e^{Lut}\,dt
=\sum_{n\ge0}w_n(\rho,p)u^n.
\label{eq:smooth-mellin-weight}
\end{equation}
Then
\begin{equation}
\boxed{
|\partial_\rho^r\partial_p^j w_n(\rho,p)|
\le L^{A_*+o(1)}\frac{(C_*L)^n}{n!}.
}
\label{eq:smooth-mellin-factorial}
\end{equation}
The same conclusion holds for a moving endpoint
$\int_0^\rho a_{\rho,p,T}(t)e^{Lut}\,dt$, and the class is stable under every
fixed finite product of such weights.
\end{lemma}

\begin{proof}
Expanding the exponential and using Fubini gives
\[
w_n(\rho,p)=\frac{L^n}{n!}
\int_0^{C_*}a_{\rho,p,T}(t)t^n\,dt,
\]
which proves \eqref{eq:smooth-mellin-factorial}. The one-dimensional
$W^{1,1}$ inequality on the fixed interval controls every endpoint value by the
corresponding $L^1$ norm together with one $t$-derivative. Repeated
$\rho$-differentiation of a moving endpoint therefore produces only finitely
many controlled endpoint jets, beginning with
\[
a_{\rho,p,T}(\rho)\frac{(L\rho)^n}{n!},
\]
and integrals of the same form. Finally, products correspond to convolution
of the compactly supported $t$-amplitudes,
so the support remains bounded and the same factorial estimate holds with a
larger fixed $C_*$. \end{proof}

\begin{lemma}[Exact HLP levels preserve the packet radial variable]
\label{lem:hlp-radial-variable}
On the absolute-convergence right edge, multiplication of finitely many
Dirichlet series changes only the coefficient of the total Dirichlet index
$M$. After insertion of the physical packet lag and the outer-denominator
dilation, the corresponding one-$\chi$ term has saddle and phase
\[
t_*=2\pi My,\qquad e(-My).
\]
Thus exact HLP levels do not create an additional packet-support coordinate;
their packet/common-prefix weights belong to the class of
Lemma~\ref{lem:smooth-mellin-class}.
\end{lemma}

\begin{proof}
If $A_1(s)\cdots A_r(s)=\sum_Mc_MM^{-s}$, then after the exact logarithmic
translation the $s$-dependence is $(Me^{-u})^{-s}$ up to a scalar amplitude.
The one-$\chi$ stationary calculation therefore depends on the single total
index $M$, giving the displayed saddle and phase. A translated factor such as
$P(s-2/L)$ merely multiplies its coefficient at $M$ by $M^{2/L}$ and does not
change the stationary variable. \end{proof}

\begin{lemma}[Cluster contour and levelwise residues]
\label{lem:cluster-levelwise}
Fix a sufficiently large constant $C_*>4$. Since the completed
three-$Z$/Cauchy construction has only finitely many normalized translated
shifts, $C_*$ may be chosen once so that every local translated pole in
$\Sigma_{{\rm tr},k}$, uniformly in $k$, lies inside
\[
\Gamma_{L}:=\{u:|u|=C_*/L\}.
\]
For all large $T$, $\Gamma_L$ also encloses $u=0$, the algebraic locator
$u_*(F_b)$ of Lemma~\ref{lem:principal-pole-cluster}, and the model point
$2/L$, while
\begin{equation}
\sup_{u\in\Gamma_L}
\left|\frac{P(1+u)}{F_b}\right|\le q<1
\label{eq:cluster-geometric-ratio}
\end{equation}
with a fixed $q$. Hence, for every packet Mellin weight $W$ analytic inside
$\Gamma_L$, the boundary integral of the bare resummed quotient equals the
absolutely convergent sum of the geometric-level contour integrals, each of
which has only the pole at $u=0$. After the fixed completed Cauchy/shift factors
are adjoined, termwise contour integration remains valid, but any additional
enclosed translated poles are retained explicitly. For $k\le K_0$ one has
\begin{equation}
\frac{1}{2\pi i}\int_{\Gamma_L}W(u)\mathcal F_{k,L}(1+u;p)\,du
=\Res_{u=0}\bigl(W\mathcal F_{k,L}\bigr)
+\sum_{u_0\in\Sigma_{{\rm tr},k}}
\Res_{u=u_0}\bigl(W\mathcal F_{k,L}\bigr).
\label{eq:completed-level-contour-residues}
\end{equation}
The holomorphic complement contributes zero. Exact local projectors are used
only for $k\le K_0$, while the full infinite leading sum is used only on the
model side. No resummed exact quotient is transported from the original Perron
line through $u_*$.
\end{lemma}

\begin{proof}
Uniformly for $|u|\le r_0$,
$P(1+u)=u^{-1}+O(1)$, while $F_b=L/2+O(1)$. Thus on $\Gamma_L$,
\[
|P(1+u)/F_b|\le 2/C_*+o(1),
\]
which proves \eqref{eq:cluster-geometric-ratio} after increasing $C_*$. The
algebraic locator $u_*=2/L+O(L^{-2})$ lies inside this circle. On $\Gamma_L$
the identity
\[
\frac{Q}{F_b-P}
=\sum_{j\ge0}\frac{QP^j}{F_b^{j+1}}
\]
therefore converges uniformly. Multiplication by the analytic $W$ preserves
uniform convergence, so contour integration may be interchanged with the sum.
Every bare geometric term on the right is meromorphic inside $\Gamma_L$ only
at $u=0$. Cauchy's theorem therefore identifies the boundary integral of the
bare resummed quotient with the sum of these complete levelwise residues. After
the fixed completed factors are restored, uniform convergence on $\Gamma_L$
still permits termwise integration. Cauchy's theorem then adds exactly the
residues of the finitely many enclosed translated poles, giving
\eqref{eq:completed-level-contour-residues}; the remaining part is holomorphic
and integrates to zero. This equality is used only after the safe levelwise
deformation and is not used to transport the resummed quotient from the
original Perron line. The model cluster $(0,2/L)$ is enclosed by the same
circle, so the resummed leading family is represented on the same local
contour. \end{proof}

\begin{lemma}[Leading HLP jet resums to the translated model cluster]
\label{lem:leading-jet-resummation}
At the principal Laurent level
\[
P(1+u)\rightsquigarrow u^{-1},
\qquad Q(1+u)\rightsquigarrow u^{-2},
\]
the complete geometric hierarchy in $Q/(F_b-P)$ resums exactly as
\begin{equation}
\boxed{
\sum_{k\ge1}F_b^{-k}u^{-k-1}
=-\frac1u+\frac1{u-F_b^{-1}}.
}
\label{eq:leading-geometric-resummation}
\end{equation}
Consequently the \emph{resummed leading Laurent family} of the completed
source has generating principal part
\begin{equation}
\boxed{-\frac2u+\frac1{u-F_b^{-1}}.}
\label{eq:leading-completed-cluster}
\end{equation}
Since $F_b=L/2+O(1)$,
\[
F_b^{-1}=\frac2L+O(L^{-2}),
\]
so \eqref{eq:leading-completed-cluster} is precisely the finite-$T$ precursor
of the model principal parts
$-2/u+1/(u-2/L)$ in \eqref{eq:Hmod}. The $O(L^{-2})$ displacement of the
translated pole is one-log lower after the common $L^2$ Perron normalization
and is included in the exact--model local comparison below.
\end{lemma}

\begin{proof}
The geometric expansion on $\Gamma_L$ is
\[
\frac{Q}{F_b-P}=\sum_{k\ge1}F_b^{-k}QP^{k-1}.
\]
Replacing only the common principal Laurent factors gives the series on the
left of \eqref{eq:leading-geometric-resummation}. Summing it as a geometric
series gives
\[
\frac1{F_bu^2}\frac1{1-(F_bu)^{-1}}
=\frac1{u(F_bu-1)}
=-\frac1u+\frac1{u-F_b^{-1}}.
\]
Adding the principal part $-u^{-1}$ of the separate $-P$ source proves
\eqref{eq:leading-completed-cluster}. Multiplication by the common shifted-zeta,
Cauchy, curvature, and packet factors does not alter this source-level pole
identity.
\end{proof}

Now write the independent Perron variable as $z=1+u$. Put
\[
A(u):=uP(1+u),\qquad
B(u):=u^2Q(1+u),\qquad
C(u):=u^2\zeta(1+u)^2.
\]
On a fixed disk $|u|\le r_0$, the Laurent expansions at the zeta pole show that
$A,B,C$ are analytic, $A(0)=B(0)=C(0)=1$, and all fixed derivatives are
bounded. For $k\ge1$ and $0\le q\le2$, differentiate the completed level $q$
times with respect to the normalized curvature shift $p$ and then set $p=0$.
A derivative falling on a factor shifted by $a=p/L$ contributes one factor
$L^{-1}$ and can raise the local pole order by one; derivatives falling on the
normalized archimedean or packet factors do not raise the pole order and need
not carry this factor. In particular the twice pole-raising contribution is
fixed by the exact curvature identity \eqref{eq:curvature-source}:
\begin{equation}
\boxed{
L^{-2}Q(1+u)^2P(1+u)^{k-1}\zeta(1+u)^2
=L^{-2}u^{-k-5}A(u)^{k-1}B(u)^2C(u).
}
\label{eq:uniform-level-factorization}
\end{equation}
The fixed shift residues commute with these derivatives by
\eqref{eq:residue-commute}. After the moving shift divisors have been canceled,
Leibniz' rule and Cauchy's estimate therefore give, for $k\le K_0$,
\begin{equation}
\boxed{
\left.\partial_p^q\mathcal F_{k,L}^{\rm ex}(1+u;p)\right|_{p=0}
=\sum_{\ell=0}^qL^{-\ell}u^{-k-3-\ell}
\bigl(H_{k,q,\ell,0}+u\Theta_{k,q,\ell,L}(u)\bigr),
\qquad 0\le q\le2.
}
\label{eq:level-Laurent}
\end{equation}
Here $\ell$ records the number of derivatives that actually raise the local
pole order. For every fixed $a$ used below,
\begin{equation}
\boxed{
|H_{k,q,\ell,0}|+|\partial_u^a\Theta_{k,q,\ell,L}(u)|
\ll L^{o(1)}C_H^k(1+k)^C,
\qquad 0\le\ell\le q\le2.
}
\label{eq:uniform-Hk}
\end{equation}
The separate ordinary face $k=0$ has fixed pole order and is covered by the
same finite estimate after enlarging the constant. The levels $k>K_0$ are not
locally projected, by Lemma~\ref{lem:HLP-high-edge}. Hence the constants in the
local Laurent error are uniform through $k\le L^{1/8}$; no fixed-$k$ asymptotic
is being summed. Package the three required curvature coefficients into the
order-two intermediate leading Taylor jet
\begin{equation}
\boxed{
\mathcal F_{k,L}^{\rm lead,(2)}(1+u;p)
:=\sum_{q=0}^2\frac{p^q}{q!}
\sum_{\ell=0}^qL^{-\ell}u^{-k-3-\ell}H_{k,q,\ell,0}.
}
\label{eq:intermediate-leading-jet}
\end{equation}
Its derivatives through order two at $p=0$ are exactly the leading terms in
\eqref{eq:level-Laurent}. This jet is not resummed from arbitrary
$k$-dependent coefficients: it is obtained by taking the same finite
$p$-derivatives of the absolutely convergent leading geometric family in
Lemma~\ref{lem:leading-jet-resummation}. Hence differentiation and geometric
resummation commute, and the translated-pole parameter remains the one given
there. No higher $p$-jet is used in the proof.
The level weight obeys
\[
\omega_{k,L}\ll C_0^kL^{-k}.
\]
For $k\le L^{1/8}$, write
\[
W_{L,\rho,p}(u)=\sum_{n\ge0}w_n(\rho,p)u^n.
\]
The finite-interval Laplace representation gives, for $r,j\le2$,
\begin{equation}
|\partial_\rho^r\partial_p^j w_n(\rho,p)|
\le C L^nB^n\frac{(1+n)^C}{n!}.
\label{eq:packet-weight-bound}
\end{equation}
Define the packet-local residue before curvature extraction by
\[
\mathcal L_{k,L}^{\rm pkt}(\rho,p)
:=\omega_{k,L}\Res_{u=0}
\bigl[\mathcal F_{k,L}(1+u;p)W_{L,\rho,p}(u)\bigr].
\]
All $p$-derivatives in the estimates below are evaluated at $p=0$. For a term
in $\partial_p^q\mathcal L_{k,L}^{\rm pkt}$ with $q_1$ derivatives falling on
$\mathcal F_{k,L}$ and $q-q_1$ on $W$, decompose the first factor according to
\eqref{eq:level-Laurent}. Its error part is a finite sum over
$0\le\ell\le q_1$ of
\[
\omega_{k,L}L^{-\ell}\Res_{u=0}
\left[u^{-k-2-\ell}\Theta_{k,q_1,\ell,L}(u)
\,\partial_p^{q-q_1}W_{L,\rho,p}(u)\big|_{p=0}\right].
\]
Expanding $\Theta_{k,q_1,\ell,L}=\sum_j\theta_{k,q_1,\ell,j}u^j$ and using
\eqref{eq:uniform-Hk} and \eqref{eq:packet-weight-bound} gives, for each
$\ell$,
\[
\ll L^{-k-\ell+o(1)}C^k(1+k)^C
\sum_{j=0}^{k+1+\ell}r_0^{-j}
\frac{(BL)^{k+1+\ell-j}(1+k)^C}
{(k+1+\ell-j)!}.
\]
Because $\ell\le q_1\le2$ is fixed and $k/L=o(1)$, the same factorial
comparison as above makes the sum uniformly bounded after factoring its $j=0$
term. Hence, uniformly for $r,q\le2$,
\begin{equation}
\boxed{
|\partial_\rho^r\partial_p^q\Delta\mathcal L_{k,L}^{\rm pkt}(\rho,0)|
\ll L^{1+o(1)}
\frac{C_1^k(1+k)^{C_1}}{(k+1)!}.
}
\label{eq:level-local-error}
\end{equation}
Only the levels $k\le K_0$ enter this local comparison. By
Lemma~\ref{lem:HLP-high-edge}, the higher exact levels retain their direct
one-$\chi$ representation: their low sidebands belong to
$E_T^{\rm HLP,hi}$ and their high sidebands belong to regular class~(vii).

The normalization is now explicit. For the component indexed by $\ell$, the
leading jet has pole order $k+3+\ell$, so residue extraction uses
$w_{k+2+\ell}$. Its natural size is
\[
L^{-k-\ell}\frac{L^{k+2+\ell}}{(k+2+\ell)!}
=L^2\frac1{(k+2+\ell)!}.
\]
Thus every increase of the local pole order is paired with exactly the same
power of $L^{-1}$ from the normalized shift. In particular, the twice
pole-raising curvature contribution in \eqref{eq:uniform-level-factorization}
uses $w_{k+4}$ together with $L^{-2}$ and still has the common $L^2$ Perron
scale. Let $K_{\rm lead,\le K_0}^{\rm pkt}$ be the sum of these leading jets for
$k\le K_0$. After the common positive Perron amplitude $L^2$ is divided out,
\eqref{eq:level-local-error} gives
\begin{equation}
\partial_\rho^r\partial_p^q
\bigl(K_{\rm ex,\le K_0}^{\rm pkt}-K_{\rm lead,\le K_0}^{\rm pkt}\bigr)
=O(L^{-1+o(1)}),
\qquad r,q\le2,
\label{eq:local-leading-error}
\end{equation}
Here and in all derivative estimates through
\eqref{eq:local-model-error}, the $p$-derivatives are evaluated at $p=0$. Let
$K_{\rm lead,\infty}^{\rm pkt}$ denote the full leading Laurent family before
freezing the translated pole. The same calculation, together with
\eqref{eq:packet-weight-bound} and \eqref{eq:uniform-Hk}, gives
\[
\left|\partial_\rho^r\partial_p^qK_{{\rm lead},k}^{\rm pkt}\right|
\ll \frac{C^k(1+k)^C}{(k+2)!},
\qquad r,q\le2,
\]
after the common $L^2$ normalization.
Therefore, for every fixed $A>0$,
\begin{equation}
\boxed{
\partial_\rho^r\partial_p^q
\bigl(K_{\rm lead,\le K_0}^{\rm pkt}-K_{\rm lead,\infty}^{\rm pkt}\bigr)
=O_A(L^{-A}),
\qquad r,q\le2.
}
\label{eq:model-truncation-tail}
\end{equation}
This is a model-side radial error; no fast one-$\chi$ carrier remains in it.
It remains to compare the full leading family with the fixed translated model. Put
\[
c_b:=\frac{L}{F_b}=\frac{2}{\mu_b}=2+O(L^{-1}).
\]
Lemma~\ref{lem:leading-jet-resummation} shows that the translated pole of the
intermediate leading family is at $u=c_b/L$. After the common $N$-prefix has
been fixed, the resummed leading residue functional is a finite combination of
bounded-support Mellin weights and fixed Cauchy factors whose dependence on
$c_b$ is analytic for $c_b$ in a fixed neighbourhood of $2$; the pole remains a
fixed positive normalized distance from the outer local contour. Lemma
\ref{lem:smooth-mellin-class} and Cauchy's estimate therefore give a uniform
bound for the $c_b$-derivative after up to two $(\rho,p)$ derivatives. The
mean-value theorem gives
\begin{equation}
\partial_\rho^r\partial_p^q
\bigl(K_{\rm lead,\infty}^{\rm pkt}-K_{\rm mod,loc}^{\rm pkt}\bigr)
=O(L^{-1}),
\qquad r,q\le2,
\label{eq:translated-pole-displacement-error}
\end{equation}
where $K_{\rm mod,loc}^{\rm pkt}$ uses the fixed model parameter $c=2$, i.e.
the pole $u=2/L$. The three differences telescope exactly:
\begin{align*}
K_{\rm ex,\le K_0}^{\rm pkt}-K_{\rm mod,loc}^{\rm pkt}
={}&(K_{\rm ex,\le K_0}^{\rm pkt}-K_{\rm lead,\le K_0}^{\rm pkt})\\
&+(K_{\rm lead,\le K_0}^{\rm pkt}-K_{\rm lead,\infty}^{\rm pkt})\\
&+(K_{\rm lead,\infty}^{\rm pkt}-K_{\rm mod,loc}^{\rm pkt}).
\end{align*}
Hence \eqref{eq:local-leading-error}, \eqref{eq:model-truncation-tail}, and
\eqref{eq:translated-pole-displacement-error} give
\begin{equation}
\boxed{
\partial_\rho^r\partial_p^q
\bigl(K_{\rm ex,\le K_0}^{\rm pkt}-K_{\rm mod,loc}^{\rm pkt}\bigr)
=O(L^{-1+o(1)}),
\qquad r,q\le2.
}
\label{eq:local-model-error}
\end{equation}

By Proposition~\ref{prop:trivial-outer-twist}, the invariant source has no
nontrivial arithmetic whitening. Consequently the exact/model local discrepancy
has normalized diagonal jump $O(L^{-1+o(1)})=o(1)$, and its fixed mixed lag
derivatives are bounded by a fixed power of $L$ on every retained shell.
The shell-local primitive criterion of Lemma~\ref{lem:two-variable-primitive},
together with the metric comparison
\eqref{eq:local-global-Gram-comparison} and the choice
$A_0>C_{\rm reg}+3$ in \eqref{eq:acyclic-A0-choice}, therefore gives
\begin{equation}
\boxed{
\|\Gfr_T^{-1/2}R_T^{\rm loc}\Gfr_T^{-1/2}\|_{S_2}^2=o(N_T).
}
\label{eq:local-HS}
\end{equation}

\begin{definition}[Finite-$T$ principal operator]
\label{def:PT}
Let $W_T$ be the stable logarithmic core length fixed in
\eqref{eq:stable-log-core}, and let $C_T$ denote restriction of the straightened
physical logarithmic coordinate to $0\le u\le W_T$. After this restriction use
the normalized lag coordinate $r=u/L$ and the unitary output rescaling
$g(u)\mapsto L^{1/2}g(Lr)$. Define $P_T$ to be the pullback, through the exact
packet and diagonal maps, this core restriction, and this rescaling, of the
\emph{model} principal-part projector
$\Pi_{\Sigma_0^{\rm mod}}$ on the pair-independent common carrier core. By
Proposition~\ref{prop:trivial-outer-twist}, the auxiliary arithmetic map
specializes to its $d=1$ coordinate and every outer-ratio translation is the
identity for the invariant contour source. Thus $P_T$ retains the finite-$T$
parameters $\lambda(\tau)$ and $\mu_b$, the common-core endpoint, and the actual
packet/common-prefix weight of
\eqref{eq:joint-weight} on $0\le u\le W_T$. No stationary Fourier-mode
restriction is imposed inside this core. The exact levels $k>K_0$ are not part
of $P_T$ and retain the direct sideband decomposition of
Lemma~\ref{lem:HLP-high-edge}; only their model leading tail is present through
the full resummed projector. The complementary logarithmic tail $W_T<u\le L$
is assigned by Proposition~\ref{prop:terminal-log-tail}. The auxiliary outward
arithmetic annulus of Proposition~\ref{prop:outer-excess-edge} is empty in the
invariant-source specialization. Slow-parameter freezing and the low-level
exact/model discrepancy are excluded from $P_T$.
\end{definition}

\begin{proposition}[Uniform-twisted packet source decomposition]
\label{prop:uniform-twisted-replacement}
After the $f'$ source, horizontal connectors, and nonstationary sector have
been separated, restrict the stationary exact HLP/Hardy source to the stable
logarithmic band $0\le u\le W_T$. Here ``stable'' refers only to the logarithmic
packet coordinate. By Proposition~\ref{prop:trivial-outer-twist} the invariant
source has no outward arithmetic annulus. On this stable band the same packet contour
matrix has the disjoint decomposition
\begin{equation}
\boxed{
A_{T}^{\rm stat,ex,stable}=P_T+E_T^{\onechi}+R_T^{\rm src}.
}
\label{eq:uniform-twisted-replacement}
\end{equation}
Here $E_T^{\onechi}$ is the direct physical one-$\chi$ family formed by the
low-sideband parts of the high exact HLP levels, the additional translated
principal parts, and the remaining direct stationary sources. The auxiliary
outer-excess family is zero for the invariant contour source.
Its rank is bounded by Proposition~\ref{prop:edge-rank}. Every complementary
high-sideband part belongs instead to $R_T^{\rm src}$, together with the other
mutually exclusive regular source classes identified in the proof below. Their relative Hilbert--Schmidt bounds are established in
Proposition~\ref{prop:regular-ledger}.
\end{proposition}

\begin{proof}
Work first on the absolute-convergence right edge and separate the horizontal
connectors and the nonstationary sector. The packet/common-prefix weight is kept
exactly at this stage; no replacement by the flat plateau is made in the present
source decomposition.

Expand the pointwise resolvent \eqref{eq:Hex-variable} geometrically on the safe
right edge before any freezing. At level $k$, multiplication of Dirichlet series
only changes the total coefficient index $M$, and Section~\ref{sec:stationary}
gives
\[
d_{k,M}(p)\,a_{k,M,T}(y;p)e(-My),
\qquad t_*=2\pi My.
\]
Split the hierarchy at $K_0=L^{1/8}$. The levels $k>K_0$ remain in this original
pointwise-$f(t)$ representation. Lemma~\ref{lem:HLP-high-edge} assigns their low
sidebands to $E_T^{\onechi}$ and their high sidebands to the direct regular
source later estimated in class~(vii).

For the finite hierarchy $k\le K_0$, insert the smooth height partition. On each
block Lemma~\ref{lem:block-freezing} replaces the pointwise source by the frozen
source \eqref{eq:Hex}. The partition has bounded overlap, so the freezing
contribution is
\[
O\!\left(L^{-D-1+C_{\rm fr}}\right)=o(1)
\]
by \eqref{eq:height-block-choice}; it belongs only to $R_T^{\rm src}$. It remains
to analyze the frozen low-level source on one block.

For each $k\le K_0$, apply the canonical decomposition
\eqref{eq:canonical-local-decomp} on the closed contour $\Gamma_L$ of
Lemma~\ref{lem:cluster-levelwise}. The exact residue identity
\eqref{eq:completed-level-contour-residues} pairs the complete Laurent operator
$\Pi_0^{(k)}$ with the actual joint packet weight of
\eqref{eq:joint-weight} and retains every enclosed translated projector
$\Pi_{\Sigma_{{\rm tr},k}}$ separately. Each translated term remains a direct
stationary one-$\chi$ source and is expanded in unit-shell Fourier sidebands:
its low sidebands belong to the edge term and its high sidebands to regular
class~(vii). After these principal parts have been subtracted, the remaining function
$\mathcal F_{k,L}^{\rm hol}$ is holomorphic throughout the interior of
$\Gamma_L$. Hence, for every analytic packet weight $W$,
\begin{equation}
\boxed{
\frac{1}{2\pi i}\int_{\Gamma_L}
W(u)\mathcal F_{k,L}^{\rm hol}(1+u;p)\,du=0.
}
\label{eq:local-holomorphic-zero}
\end{equation}
Thus no independent ``near-pole analytic remainder'' survives the local closed
contour. By \eqref{eq:residue-commute}, the same levelwise extraction applies
to all three completed $Z$-faces, and
Lemma~\ref{lem:smooth-mellin-class} gives the required factorial Mellin bounds.
For $k\le K_0$, the exact stationary level has the mutually exclusive form
\[
\mathcal T_{k}^{\rm stat}
=\mathcal L_{k}^{\rm ex}
+\mathcal E_{k}^{\rm tr}
+\mathcal D_{k}^{\rm dir}.
\]
Here $\mathcal L_k^{\rm ex}$ is the complete common local Laurent contribution,
$\mathcal E_k^{\rm tr}$ the additional translated principal parts, and
$\mathcal D_k^{\rm dir}$ the complementary pole-free right-edge one-$\chi$
contribution. The levels $k>K_0$ are not subjected to this local projector and
remain in the direct edge representation of Lemma~\ref{lem:HLP-high-edge}. Any smooth transition used to join the
local contour to the complementary Mellin contour is supported away from the
pole cluster and is included in $\mathcal D_k^{\rm dir}$.

After removing the unprojected levels $k>K_0$ by
Lemma~\ref{lem:HLP-high-edge}, apply the trivial HLZ outer-twist specialization
of Proposition~\ref{prop:trivial-outer-twist}. Its outer indices are all equal
to $1$ and the auxiliary divisor coordinate is the identity; this notation is
independent of the HLP level index $k$ used throughout the present section.
Accordingly the outward annulus of Proposition~\ref{prop:outer-excess-edge} is
empty. On the resulting
common carrier core, replace \eqref{eq:level-Laurent} first by its intermediate
leading jet and then replace the resummed pole parameter $c_b=L/F_b$ by the
fixed model value $2$. The resulting principal family, with the actual
finite-$T$ packet weight, defines $P_T$. Proposition~\ref{prop:X-tagging}
identifies the explicit $X$-powers with $\lambda(\tau)$, while the reassembled
self term retains $\mu_b=2F_b/L$. No nontrivial outer translation or
outer-arithmetic factor occurs in the invariant contour source.

The local regular term $R_T^{\rm loc}$ is the sum of the low-level exact/leading
discrepancy \eqref{eq:local-leading-error}, the model truncation tail
\eqref{eq:model-truncation-tail}, and the translated-pole displacement
\eqref{eq:translated-pole-displacement-error}; their combined estimate is
\eqref{eq:local-HS}. Only in
Section~\ref{sec:model-floor}, for the purpose of proving a Loewner floor, are
the remaining slow parameters of $P_T$ compared with the frozen reference
model $A_T^G$.

The local holomorphic remainder contributes zero by
\eqref{eq:local-holomorphic-zero}. On the complementary Mellin region, the
original absolutely convergent Dirichlet expansion is retained term by term,
so $\mathcal D_k^{\rm dir}$ keeps the physical carrier $e(-My)$. Its pole-free
transition factors remain in the smooth stationary-symbol class of
Lemma~\ref{lem:stationary-symbol}. Expanding that unit-shell symbol in Fourier
sidebands, the low sidebands lie in the common macroscopic shell-edge range and
hence have rank $o(N_T)$, while the high sidebands satisfy part (vii) of
Proposition~\ref{prop:regular-ledger}. No diagonal-jump estimate for the
order-one complementary source is used.

Finally Lemma~\ref{lem:HLP-high-edge} sends the low sidebands of $k>K_0$ to the
edge matrix and their high sidebands to regular class~(vii). The local estimates
are summed only over $k\le K_0$, with the model truncation tail
\eqref{eq:model-truncation-tail}. This gives
\eqref{eq:uniform-twisted-replacement}. The direct carrier structure of the
complementary source is inherited from the original levelwise right-edge
Dirichlet expansion; its quantitative high-sideband estimate is deferred to
Proposition~\ref{prop:regular-ledger}.
\end{proof}
